\chapter{n 维空间.}

我们已经有机会说过，平面和空间的解析几何的发展如何引导数学家引入 n 维空间的概念；这个概念为他们提供了一种几何语言，极其便于以简洁明了的方式表达关于任意多个变量的方程的代数定理，尤其是线性代数的一切一般结果（见第 61 页）。但是，尽管到十九世纪中叶这种语言已在众多几何学家中间通行，它仍然是纯粹约定性的，而且由于缺乏三维以上空间的“直观”表象，在这些空间中，在平面或空间里仅凭“直觉”便被允许的那种“按连续性”的论证似乎遭到了禁止。黎曼在其关于位置分析和几何学基础的研究中，第一个敢于通过与三维空间的类比按这种方式进行论证（见第 176 页）；\(^{1}\)在他之后，众多数学家开始卓有成效地使用这类论证，尤其是在多个复变量的代数函数理论中。但由于当时对空间直观的把握非常有限，人们有理由对这类思考在证明中的价值保持怀疑，只承认它们纯粹是启发性的，因为它们只是使某些定理的正确性显得颇为可能。因此，H. 庞加莱在其 1887 年关于两个复变量的函数的二重积分的留数的论文中，尽可能避免一切对四维空间直观的求助：“由于这种超几何语言至今仍令许多优秀的头脑感到反感”，他说，“我只偶尔使用它”；他为达到这一目的而使用的种种“技巧”，使他得以回到三维空间中的拓扑论证，在那里他也就不再犹豫去诉诸直观了（{[}251 a{]}，第 III 卷，第 443 页及以下）。

此外，康托尔的种种发现，尤其是确立 R 与 \(R^{n}\) 等势的那条著名定理（它似乎使维数概念本身也成了问题）\(^{2}\)表明，为了给几何学与拓扑学的论证提供坚实的基础，必须使它们完全摆脱对直观的任何求助。我们已经说过（参见第 139 页），这一需要正是一般拓扑的现代观念的起源；但甚至在这一学科创立之前，对数空间的拓扑及其最直接的推广（“n 维流形”）的严格研究就已经开始了，其所用方法主要源自拓扑学中称为“组合拓扑”或者更好地说“代数拓扑”的那个分支，关于它我们在这里无法论述。
% semantic-fragments: legacy-input-seam
\relax{}
